% ===== latex-recomposed:document =====
% flm: style file flmrecomposed.c57af267.sty
\begin{defterm}{Entropy}\phantomsection\flmLabelsWithDisplayText{Entropy}{\flmLabels{\label{ref1}\label{ref2}}}\ignorespaces
A measure.
\end{defterm}

See \flmRef{ref2}.
%%
% ===== latex-recomposed:style flmrecomposed.c57af267.sty =====
\NeedsTeXFormat{LaTeX2e}
\ProvidesPackage{flmrecomposed.c57af267}
% --- packages ---
% flm.latex.pkg.hyperref
\RequirePackage{hyperref}
% flm.latex.pkg.zref-clever
\RequirePackage{zref-clever}
% --- commands and environments the library defines for the document ---
% flm.latex.def.defterm
% Defined terms. A term definition is written as the environment defterm,
% whose one argument is the term. This file defines the environment further
% below, in the part headed "definitions written from the configuration of
% FLM features"; the environment's opening text is the lead-in that the
% document was recomposed with. The body of a term definition opens with
% \phantomsection, so that a link to the term definition leads to it, and
% with \flmLabelsWithDisplayText, so that a reference to one of its labels
% prints the term.
% \flmTerm{<label>}{<shown>}: <shown>, a link to the definition of a term,
%   <label> being a label of that definition, as \flmHyperref takes it.
% \flmTermInDefterm{<shown>}: <shown> written inside the definition of its
%   own term: the term being defined, set apart, with no link.
% Redefine either command to set term links another way.
\ProvideDocumentCommand\flmTerm{m m}{\flmHyperref{#1}{#2}}
\ProvideDocumentCommand\flmTermInDefterm{m}{\emph{#1}}
% What a term link contributes to a PDF bookmark: its text alone. The code
% below runs when the document begins (\AtBeginDocument) and does nothing
% where hyperref is not loaded.
\AtBeginDocument{\ifdefined\pdfstringdefDisableCommands
  \pdfstringdefDisableCommands{%
    \def\flmTerm#1#2{#2}%
    \def\flmTermInDefterm#1{#1}%
  }\fi}
% flm.latex.def.refs
% \flmRef[<options>]{<label>}: a reference that gives the name of its target
% ("Section 2"), set by \zcref of the zref-clever package.
%
% The options are the ones the document wrote, in FLM's own vocabulary. This
% definition translates each of them by its key -- a value written after the
% key is ignored, as in FLM -- and drops every key it has no translation for:
%
%   noprefix    -> noname                 the number or title alone, with no
%                                         name before it ("2", not
%                                         "Section 2")
%   S           -> S                      the first name capitalized and
%                                         unabbreviated, for a reference that
%                                         opens a sentence
%   capital     -> cap=true,abbrev=false  the names capitalized ("Lemma 1")
%   lowercase   -> cap=false,abbrev=false the names in lowercase ("lemma 1")
%   abbreviated -> abbrev=true            the abbreviated names ("Lem. 1")
%                                         where the reference type declares
%                                         them, the full names otherwise
%   anything else -> dropped              the default wording; a variant of
%                                         your own naming has no counterpart
%                                         in zref-clever, which would stop on
%                                         a key it does not know
%
% Redefine \flmRef to set references differently -- another package, another
% wording, your own variants, no link. Your definition sees the options as the
% document wrote them, in FLM's vocabulary, and maps them as it likes.
%
% The translation uses expl3, LaTeX's programming layer, which is part of the
% format and needs no package; its names read only between \ExplSyntaxOn and
% \ExplSyntaxOff. \keyval_parse:NNn reads the option list as LaTeX reads any
% key-value list: it hands a key written alone to the first macro below and a
% key with its value to the second, which drops the value. \str_case:nnF looks
% the key up in the table and appends what \zcref is to be given, or nothing.
% The macros with @ in their names are internals, not a customization point:
% redefine \flmRef, not its parts.
\ExplSyntaxOn
\def\flm@ref@add#1{\tl_put_right:Nn \flm@ref@options {#1,}}
\def\flm@ref@key#1{\str_case:nnF {#1}
  {{noprefix}{\flm@ref@add{noname}}
   {S}{\flm@ref@add{S}}
   {capital}{\flm@ref@add{cap=true,abbrev=false}}
   {lowercase}{\flm@ref@add{cap=false,abbrev=false}}
   {abbreviated}{\flm@ref@add{abbrev=true}}}
  {}}
\def\flm@ref@pair#1#2{\flm@ref@key{#1}}
\def\flm@ref@zcref#1#2{\zcref[#1]{#2}}
\ProvideDocumentCommand\flmRef{O{} m}{%
  \def\flm@ref@options{}%
  \keyval_parse:NNn \flm@ref@key \flm@ref@pair {#1}%
  \exp_args:No \flm@ref@zcref {\flm@ref@options} {#2}}
\ExplSyntaxOff
% flm.latex.def.hyperref
% \flmHyperref{<label>}{<text>}: <text>, shown as a link to the place <label>
% refers to.
%
% Making a link in a PDF takes the hyperref package, and loading it changes a
% whole document -- every reference becomes a link, bookmarks are built. That
% is the author's decision, not this file's, so this file neither loads
% hyperref nor asks for it: the definition below links through it where the
% document loaded it and prints the text alone where it did not, so the
% document compiles either way. The question is about this document's own
% preamble, so the answer is the same wherever the document is compiled; it is
% asked here, at the reference, so it is right whether this file was loaded
% before hyperref or after it.
%
% Redefine \flmHyperref to link differently -- another package, a footnote
% with the target's number, no link at all.
\ProvideDocumentCommand\flmHyperref{m m}{%
  \IfPackageLoadedTF{hyperref}{\hyperref[{#1}]{#2}}{#2}}
% What such a reference contributes to a PDF bookmark: the text alone. A
% bookmark is plain text and carries no link, and without this line the label
% would appear in it beside the text. The hook runs when the document begins
% and does nothing where hyperref is absent.
\AtBeginDocument{\ifdefined\pdfstringdefDisableCommands
  \pdfstringdefDisableCommands{%
    \def\flmHyperref#1#2{#2}%
  }\fi}
% flm.latex.def.label-display-text
% \flmLabelsWithDisplayText{<text>}{<labels>}: makes <labels> refer to this
% place, with <text> as what a reference to them prints. It is the only
% command for that: a heading the class did not number uses it, and so does
% anything else that has a name but no number.
%
% It places no anchor of its own. Under hyperref a sectioning command, starred
% or not, already places one at the title; a \label after it points there.
% \refstepcounter would place a second anchor here, below the title, and a
% link would lead below the heading. So the counter is stepped with
% the plain \stepcounter and the two things this needs are set here:
% \@currentlabel, which is what a \label records, and \@currentcounter,
% which is what zref-clever reads to give the label its reference type. The
% group keeps every change local.
%
% <text> is expanded to make \@currentlabel, and \label does not survive that
% expansion; so a \label inside <text> (for example a footnote's) is dropped
% from the label's text only. At its own place in the document it still makes
% its own label.
%
% A caller with no sectioning command before it -- for example a defined term
% in running text -- writes \phantomsection at its own place, so that a link
% has an anchor of its own to lead to.
%
% Redefine \flmLabelsWithDisplayText to change how such labels are made.
\ifcsname c@flm@internal@refcnt\endcsname\else\newcounter{flm@internal@refcnt}\fi
\ProvideDocumentCommand\flmLabelsWithDisplayText{m m}{%
  \begingroup
    \stepcounter{flm@internal@refcnt}%
    {\let\label\@gobble\protected@xdef\flm@lwdt@text{#1}}%
    \let\@currentlabel\flm@lwdt@text
    \def\@currentcounter{flm@internal@refcnt}%
    #2%
  \endgroup}
% flm.latex.def.zref-clever-type
% The zref-clever reference type of the counter \flmLabelsWithDisplayText
% steps, with every name prefix empty, so that a reference to such a label
% prints its text alone ("Introduction", not "Section Introduction").
% Every other label keeps the name its own type gives it. Redefine these
% settings, or the type of the counter, to put something before the text.
\zcsetup{countertype={flm@internal@refcnt=flm--internal--ref}}
\zcRefTypeSetup{flm--internal--ref}{Name-sg = , name-sg = , Name-pl = , name-pl = ,}
% flm.latex.def.labels
% \flmLabels{<\label commands>}: the labels of one place in the document.
% In text it runs them as written. On a line of a display-math
% environment, where amsmath keeps the last label of a line and loses the
% ones before it, it runs the first and records every further one as an alias:
% a line per alias is written to the .aux file after the last page and read on
% the next run, making the alias refer to what the first label refers to
% (for \ref and for \zcref alike). Redefine \flmLabels to change what a
% group of labels does; \def\flmLabels#1{#1} restores plain LaTeX behavior,
% amsmath's "Multiple \label's" message included.
\ProvideDocumentCommand\flmLabels{m}{%
  \ifmmode\expandafter\flm@labels@math\else\expandafter\@firstofone\fi{#1}}
% Reading the group: the first \label as written, the rest as aliases of it.
\def\flm@labels@math#1{\def\flm@lbl@first{}\flm@lbl@scan#1\flm@lbl@stop}
\def\flm@lbl@scan{\@ifnextchar\label\flm@lbl@take\flm@lbl@rest}
\def\flm@lbl@take\label#1{%
  \ifx\flm@lbl@first\@empty
    \def\flm@lbl@first{#1}%
    \label{#1}%
  \else
    \expandafter\flm@recordalias\expandafter{\flm@lbl@first}{#1}%
  \fi
  \flm@lbl@scan}
\def\flm@lbl@rest#1\flm@lbl@stop{#1}
% Collecting the aliases. A name is recorded once however often the line is
% typeset (amsmath measures a line before setting it).
\def\flm@aliaslist{}
\def\flm@recordalias#1#2{%
  \ifcsname flm@alias@#2\endcsname\else
    \expandafter\gdef\csname flm@alias@#2\endcsname{}%
    \g@addto@macro\flm@aliaslist{\flm@alias@entry{#2}{#1}}%
  \fi}
% Writing them, from the hook that runs after the last page has been shipped
% out: every \newlabel of that page has reached the .aux file by then, so the
% alias lines come after them. \AtEndDocument would run too early.
\def\flm@alias@write#1#2{\immediate\write\@auxout{\string\flm@aliaslabel{#1}{#2}}}
\AddToHook{enddocument/afterlastpage}{\let\flm@alias@entry\flm@alias@write\flm@aliaslist}
% Reading them back: the alias refers to what the original refers to. The
% assignments are global because the .aux file is read inside a group.
\def\flm@aliaslabel#1#2{%
  \ifcsname r@#2\endcsname
    \global\expandafter\let\csname r@#1\expandafter\endcsname\csname r@#2\endcsname
  \fi
  \ifcsname Z@R@#2\endcsname
    \global\expandafter\let\csname Z@R@#1\expandafter\endcsname\csname Z@R@#2\endcsname
  \fi}
% --- definitions written from the configuration of FLM features ---
% defterm (the defterm feature's, written from its configuration): defined unless the name is already taken, so a definition of it made before this file is loaded wins; \RenewDocumentEnvironment after loading this file changes it
\ProvideDocumentEnvironment{defterm}{+m}{}{}
\endinput
